\documentclass[a4paper,12pt]{article}
\usepackage[utf8]{inputenc}
\usepackage[T1]{fontenc}
\usepackage[french]{babel}
\usepackage{amsmath, amssymb, amsfonts}
\usepackage{array, longtable, booktabs}
\usepackage[table]{xcolor}
\usepackage{geometry}
\geometry{left=1.8cm, right=1.8cm, top=2cm, bottom=2cm}
\usepackage{enumitem}
\usepackage{hyperref}
\hypersetup{colorlinks=true, linkcolor=blue, urlcolor=blue}

% Commandes pour les coches
\newcommand{\fait}{\ensuremath{\checkmark}}
\newcommand{\revoit}{\ensuremath{\circlearrowright}}
\newcommand{\casevide}{\ensuremath{\square}}

% Types de colonnes robustes
\newcolumntype{P}[1]{>{\raggedright\arraybackslash}p{#1}}
\newcolumntype{C}[1]{>{\centering\arraybackslash}p{#1}}

% Pour que les longtable utilisent toute la largeur
\setlength{\LTleft}{0cm}
\setlength{\LTright}{0cm}

\title{Programmation annuelle de Mathématiques -- STS BIO AC (2\up{me} année) \\
       Année 2026--2027 -- Zone C (3 h/semaine)}
\author{}
\date{Version du \today}

\begin{document}

\maketitle

\section*{Présentation générale}
Ce document propose une organisation spiralaire de l’enseignement des mathématiques en 2\up{me} année de STS Bioanalyses en laboratoire de contrôle, conforme au programme officiel et adapté aux acquis (ou lacunes) de la 1\up{re} année. L’année comporte 31 semaines de cours effectifs (hors vacances et jours fériés).

\medskip
\noindent \textbf{Prise en compte du bilan de 1\up{re} année :}
\begin{itemize}
    \item Les notions \textcolor{green}{vertes} dans le fichier de suivi sont considérées comme acquises (elles sont rapidement réactivées si nécessaire).
    \item Les notions \textcolor{red}{rouges} (ou non renseignées) feront l’objet d’un enseignement explicite ou d’un réinvestissement approfondi en début d’année, afin de consolider les prérequis indispensables à la suite du programme.
\end{itemize}

\medskip
\textbf{Place de Python :} La programmation est intégrée de façon transverse et évaluée dans \emph{tous} les sujets (formatives, contrôles, DM, PBL). Les compétences algorithmiques (boucles, tests, fonctions, simulation, calcul numérique) sont développées en parallèle des mathématiques.

\medskip
\textbf{Utilisation de l’IA dans les DM :} Dans chaque devoir maison, l’élève est invité à utiliser un outil d’IA générative (ChatGPT, Mistral, Copilot, etc.) pour :
\begin{itemize}
    \item obtenir des explications ou des exemples supplémentaires,
    \item suggérer des pistes de résolution,
    \item débuguer ou améliorer un code Python,
    \item confronter différentes méthodes.
\end{itemize}
\emph{L’IA ne doit pas rédiger la solution finale.} Un \textbf{journal de bord} retraçant les échanges avec l’IA est annexé à chaque DM.

\section{Tableau récapitulatif des contenus à enseigner}
Le tableau ci-dessous reprend l’intégralité des attendus du programme de 2\up{me} année, en précisant le statut des prérequis de 1\up{re} année (vert = déjà traité, rouge = à reprendre).

\rowcolors{2}{gray!8}{white}
\begin{longtable}{|P{9cm}|C{2.5cm}|C{2.5cm}|}
    \hline
    \textbf{Thème / Contenu} & \textbf{Acquis en 1\up{re}} & \textbf{À reprendre en 2\up{me}} \\
    \hline
    \endhead
    \hline
    \multicolumn{3}{r}{Suite de la page suivante} \\
    \endfoot
    \hline
    \endlastfoot

    % --- Rappels (1ère non traités) ---
    \multicolumn{3}{|P{9cm}|}{\textbf{Rappels et compléments (1\up{re} non traités)}} \\
    Signe et nullité d’une fraction & \casevide & \fait \\
    Simplification de fractions (décimaux / notation scientifique) & \casevide & \fait \\
    Pourcentages « parallèles » & \casevide & \fait \\
    Pourcentages « successifs » & \casevide & \fait \\
    Pourcentage de pourcentage & \casevide & \fait \\
    Variation absolue / relative (taux d’évolution) & \casevide & \fait \\
    Taux d’évolution global & \casevide & \fait \\
    Taux d’évolution moyen & \casevide & \fait \\
    Taux d’évolution réciproque & \casevide & \fait \\
    Indice simple en base 100 & \casevide & \fait \\
    Systèmes positionnels (binaire, hexadécimal, conversions) & \casevide & \fait \\
    \hline

    % --- Compléments d'analyse ---
    \multicolumn{3}{|P{9cm}|}{\textbf{Compléments d’analyse (1\up{re} non traités)}} \\
    Fonctions polynômes, racine carrée & \casevide & \fait \\
    Fonctions logarithme, sin, cos, tan & \casevide & \fait \\
    Asymptotes & \casevide & \fait \\
    Dérivée d’un quotient & \casevide & \fait \\
    Recherche d’extrema, tracé de courbes & \casevide & \fait \\
    Problèmes d’optimisation & \casevide & \fait \\
    \hline

    % --- Statistique descriptive approfondie ---
    \multicolumn{3}{|P{9cm}|}{\textbf{Statistique descriptive approfondie}} \\
    Regroupement en classes & \casevide & \fait \\
    Fréquences cumulées & \casevide & \fait \\
    Histogramme, polygone des fréquences & \casevide & \fait \\
    Diagramme en boîte (définition, lecture, construction) & \casevide & \fait \\
    Médiane, quartiles, mode & \casevide & \fait \\
    Étendue, écart-type, variance & \casevide & \fait \\
    \hline

    % --- Probabilités – compléments ---
    \multicolumn{3}{|P{9cm}|}{\textbf{Probabilités – compléments}} \\
    Indépendance de deux événements & \casevide & \fait \\
    Variance d’une variable aléatoire discrète & \casevide & \fait \\
    Approximation par Poisson & \casevide & \fait \\
    Approximation par la normale & \casevide & \fait \\
    \hline

    % --- Calcul intégral ---
    \multicolumn{3}{|P{9cm}|}{\textbf{Calcul intégral (2\up{me} année)}} \\
    Sommes de Riemann, approche graphique & \casevide & \fait \\
    Linéarité, relation de Chasles, positivité & \casevide & \fait \\
    Primitives, lien dérivation/intégration & \casevide & \fait \\
    Théorème fondamental de l’analyse & \casevide & \fait \\
    Calcul d’aires (entre deux courbes) & \casevide & \fait \\
    Changement de variable (substitution) & \casevide & \fait \\
    Calcul de volumes, valeur moyenne & \casevide & \fait \\
    Applications en cinétique & \casevide & \fait \\
    \hline

    % --- Équations différentielles ---
    \multicolumn{3}{|P{9cm}|}{\textbf{Équations différentielles}} \\
    Équations à variables séparables & \casevide & \fait \\
    Équations linéaires (solution générale, particulière) & \fait & \casevide \\
    Problème de Cauchy (condition initiale) & \fait & \casevide \\
    Modèles de croissance / décroissance exponentielle & \fait & \casevide \\
    Applications aux dosages, réactions chimiques, transferts thermiques & \casevide & \fait \\
    \hline

    % --- Probabilités 2 ---
    \multicolumn{3}{|P{9cm}|}{\textbf{Probabilités 2 (variables continues)}} \\
    Densité de probabilité, fonction de répartition & \casevide & \fait \\
    Espérance et variance (variable continue) & \casevide & \fait \\
    Paramètres de la loi normale, propriétés & \casevide & \fait \\
    Loi normale centrée réduite, table/logiciel & \casevide & \fait \\
    Calcul de probabilités, quantiles & \casevide & \fait \\
    Théorème central limite (notion) & \casevide & \fait \\
    \hline

    % --- Statistique inférentielle ---
    \multicolumn{3}{|P{9cm}|}{\textbf{Statistique inférentielle}} \\
    Distribution d’échantillonnage (moyenne, proportion) & \casevide & \fait \\
    Estimation ponctuelle (moyenne, variance, proportion) & \casevide & \fait \\
    Intervalle de confiance (moyenne $\sigma$ connu/inconnu, proportion) & \casevide & \fait \\
    Principe des tests, risques ($\alpha$, $\beta$), p-valeur & \casevide & \fait \\
    Test sur la moyenne, test sur une proportion & \casevide & \fait \\
    Applications : cartes de contrôle, conformité & \casevide & \fait \\
    \hline
\end{longtable}

\section{Calendrier et programmation harmonieuse}

\subsection{Calendrier retenu (zone C, jours ouvrés)}
Les vacances et jours fériés ont été pris en compte. Les semaines de cours sont les suivantes :

\begin{center}
    \begin{tabular}{|C{3.5cm}|C{5cm}|C{2.5cm}|}
        \hline
        \textbf{Période} & \textbf{Dates} & \textbf{Nb de semaines} \\
        \hline
        Rentrée – Toussaint & 1er sept. – 16 oct. 2026 & 7 \\
        Toussaint – Noël & 2 nov. – 18 déc. 2026 & 7 \\
        Noël – Hiver & 4 janv. – 5 fév. 2027 & 5 \\
        Hiver – Printemps & 22 fév. – 2 avr. 2027 & 6 \\
        Printemps – Fin mai & 19 avr. – 31 mai 2027 & 6 \\
        \hline
        \textbf{Total} & & \textbf{31 semaines} \\
        \hline
    \end{tabular}
\end{center}

\textit{Jours fériés exclus :} 1er janvier, 29 mars (Pâques), 6 mai (Ascension), 17 mai (Pentecôte).

\subsection{Programmation par trimestre (approche spiralée)}
Les tableaux suivants précisent pour chaque semaine les contenus, les évaluations formatives (avec Python) et les remarques.

% --- 1er trimestre ---
\newpage
\subsubsection{1\up{er} trimestre (septembre – novembre 2026) – 14 semaines}

\rowcolors{2}{gray!8}{white}
\begin{longtable}{|C{1.2cm}|P{4.0cm}|P{5.5cm}|P{2.5cm}|}
    \hline
    \textbf{Sem.} & \textbf{Contenus nouveaux} & \textbf{Évaluation formative -- points évalués (dont Python)} & \textbf{Remarques} \\
    \hline
    \endhead
    \hline
    \multicolumn{4}{r}{Suite de la page suivante} \\
    \endfoot
    \hline
    \endlastfoot

    1-2 & \textbf{Rappels – Pourcentages et évolutions.}
    Pourcentages parallèles / successifs / de pourcentage.
    Taux d’évolution (absolu, relatif, global, moyen, réciproque). Indices. & Calculs de taux. \textbf{Python} : fonction qui calcule un taux d’évolution moyen à partir d’une liste de taux. & Réactivation indispensable pour la chimie et la gestion de données. \\
    \hline
    3-4 & \textbf{Rappels – Numération et simplification.}
    Systèmes positionnels (binaire, hexadécimal). Conversions.
    Simplification de fractions (décimaux, notation scientifique). Signe. & Conversions décimal $\leftrightarrow$ binaire. \textbf{Python} : fonctions de conversion. & Utile pour les codages et les calculs d’incertitudes. \\
    \hline
    5-6 & \textbf{Compléments d’analyse.}
    Fonctions polynômes, racine carrée, logarithme, sinus, cosinus, tangente. & Dérivées et primitives de ces fonctions. \textbf{Python} : tracé de courbes avec \texttt{matplotlib}. & Lien avec les modèles en biologie (croissance, pH). \\
    \hline
    7-8 & \textbf{Compléments d’analyse – suite.}
    Asymptotes, dérivée d’un quotient. Recherche d’extrema, tracé de courbes. Problèmes d’optimisation. & \textbf{Spirale :} dérivation (sem 5-6). \textbf{Python} : recherche de maximum par balayage. & Optimisation (coûts, rendements). \\
    \hline
    9-10 & \textbf{Statistique descriptive approfondie.}
    Regroupement en classes, fréquences cumulées. Histogramme, polygone des fréquences. & \textbf{Spirale :} effectifs/fréquences (1ère). \textbf{Python} : tracé d’histogrammes avec \texttt{matplotlib}. & Représentation des données de laboratoire. \\
    \hline
    11-12 & \textbf{Statistique descriptive – suite.}
    Diagramme en boîte (définition, lecture, construction). Médiane, quartiles, mode. Étendue, écart-type, variance. & \textbf{Spirale :} moyenne (1ère). \textbf{Python} : calcul d’écart-type et tracé de box-plot. & Paramètres de dispersion essentiels pour l’inférence. \\
    \hline
    13-14 & \textbf{Probabilités – compléments.}
    Indépendance de deux événements. Variance d’une variable discrète. Approx. par Poisson et normale. & \textbf{Spirale :} conditionnelles, binomiale (1ère). \textbf{Python} : simulation de la loi de Poisson. & Préparation aux lois continues. \\
    \hline
    \multicolumn{4}{l}{} \\
    \hline
    \multicolumn{4}{l}{\textbf{Contrôles continus :} semaines 5, 10, 14 (avec exercice Python). \textbf{Formatives :} chaque semaine.} \\
    \hline
\end{longtable}

% --- 2e trimestre ---
\newpage
\subsubsection{2\up{e} trimestre (novembre 2026 – février 2027) – 12 semaines}

\rowcolors{2}{gray!8}{white}
\begin{longtable}{|C{1.2cm}|P{4.0cm}|P{5.5cm}|P{2.5cm}|}
    \hline
    \textbf{Sem.} & \textbf{Contenus nouveaux} & \textbf{Évaluation formative -- points évalués (dont Python)} & \textbf{Remarques} \\
    \hline
    \endhead
    \hline
    \multicolumn{4}{r}{Suite de la page suivante} \\
    \endfoot
    \hline
    \endlastfoot

    15-16 & \textbf{Calcul intégral – Primitives et Riemann.}
    Sommes de Riemann, approche graphique. Primitives, lien dérivation/intégration. & \textbf{Spirale :} dérivation (T1). \textbf{Python} : méthode des rectangles pour approcher une intégrale. & Lien avec l’aire sous la courbe en cinétique. \\
    \hline
    17-18 & \textbf{Calcul intégral – Propriétés.}
    Linéarité, relation de Chasles, positivité. Théorème fondamental de l’analyse. & \textbf{Spirale :} primitives (sem 15-16). \textbf{Python} : calcul exact et approché d’intégrales. & \\
    \hline
    19-20 & \textbf{Calcul intégral – Aires et changement de variable.}
    Calcul d’aires (entre deux courbes). Changement de variable (substitution). & \textbf{Spirale :} théorème fondamental (sem 17-18). \textbf{Python} : calcul d’aire par Monte-Carlo. & \textbf{PBL 1 (2h)} : Optimisation d’un volume (ex. réacteur). \\
    \hline
    21-22 & \textbf{Calcul intégral – Applications.}
    Calcul de volumes, valeur moyenne d’une fonction. Applications en cinétique (vitesse de réaction). & \textbf{Spirale :} intégration (sem 19-20). \textbf{Python} : calcul de valeur moyenne sur un intervalle. & Lien avec la détermination de quantités de matière. \\
    \hline
    23-24 & \textbf{Équations différentielles.}
    Équations à variables séparables. Révision des équations linéaires du 1er ordre (solution générale, particulière). & \textbf{Spirale :} primitives (T2). \textbf{Python} : résolution numérique par la méthode d’Euler. & \\
    \hline
    25-26 & \textbf{Équations différentielles – Applications.}
    Problème de Cauchy. Modèles de croissance/décroissance exponentielle.
    Applications aux dosages, réactions chimiques, transferts thermiques. & \textbf{Spirale :} équadiff (sem 23-24). \textbf{Python} : simulation de la décroissance radioactive. & \textbf{PBL 2 (2h)} : Modélisation cinétique d’une réaction. \\
    \hline
    \multicolumn{4}{l}{} \\
    \hline
    \multicolumn{4}{l}{\textbf{Contrôles continus :} semaines 18, 22, 26 (avec exercice Python). \textbf{Formatives :} chaque semaine.} \\
    \hline
\end{longtable}

% --- 3e trimestre ---
\newpage
\subsubsection{3\up{e} trimestre (février – mai 2027) – 12 semaines}

\rowcolors{2}{gray!8}{white}
\begin{longtable}{|C{1.2cm}|P{4.0cm}|P{5.5cm}|P{2.5cm}|}
    \hline
    \textbf{Sem.} & \textbf{Contenus nouveaux} & \textbf{Évaluation formative -- points évalués (dont Python)} & \textbf{Remarques} \\
    \hline
    \endhead
    \hline
    \multicolumn{4}{r}{Suite de la page suivante} \\
    \endfoot
    \hline
    \endlastfoot

    27-28 & \textbf{Probabilités 2 – Variables continues.}
    Densité de probabilité, fonction de répartition. Espérance et variance (variable continue). & \textbf{Spirale :} variables discrètes (T1). \textbf{Python} : simulation de variables à densité. & \\
    \hline
    29-30 & \textbf{Loi normale.}
    Paramètres, propriétés. Loi normale centrée réduite. Utilisation de la table (ou logiciel). Calcul de probabilités et quantiles. & \textbf{Spirale :} densités (sem 27-28). \textbf{Python} : calcul de probabilités avec \texttt{scipy.stats}. & Essentiel pour le contrôle qualité. \\
    \hline
    31-32 & \textbf{Théorème central limite. Échantillonnage.}
    Notion de TCL. Distribution d’échantillonnage de la moyenne et de la proportion. & \textbf{Spirale :} loi normale (sem 29-30). \textbf{Python} : simulation de la distribution des moyennes. & Préparation à l’inférence. \\
    \hline
    33-34 & \textbf{Estimation et intervalles de confiance.}
    Estimation ponctuelle (moyenne, variance, proportion).
    IC pour la moyenne ($\sigma$ connu / inconnu). IC pour une proportion. & \textbf{Spirale :} échantillonnage (sem 31-32). \textbf{Python} : calcul d’IC à partir d’un échantillon. & \\
    \hline
    35-36 & \textbf{Tests d’hypothèses.}
    Principe, risques ($\alpha$, $\beta$), p-valeur. Test sur la moyenne, test sur une proportion.
    Applications : cartes de contrôle, conformité. & \textbf{Spirale :} estimation (sem 33-34). \textbf{Python} : implémentation d’un test de Student. & \textbf{PBL 3 (2h)} : Contrôle qualité par test d’hypothèse. \\
    \hline
    37-38 & \textbf{Synthèse et révisions.} & \textbf{Bilan général :} tous les thèmes + Python. & Préparation aux épreuves. \\
    \hline
    \multicolumn{4}{l}{} \\
    \hline
    \multicolumn{4}{l}{\textbf{Contrôles continus :} semaines 30, 34, 38 (avec exercice Python). \textbf{Formatives :} chaque semaine.} \\
    \hline
\end{longtable}

\section{Évaluations}

\subsection{Évaluations formatives (hebdomadaires)}
Chaque semaine, une courte activité (5–10 min) mêle mathématiques et Python. Les points évalués sont détaillés dans les tableaux ci-dessus.

\subsection{Évaluations sommatives (contrôles continus)}
Trois contrôles par trimestre, d’une durée de 1 h chacun. Chaque contrôle comporte une partie sans calculatrice et une partie avec Python.

\begin{center}
    \begin{tabular}{|P{3.5cm}|P{3.5cm}|P{3.5cm}|P{3.5cm}|}
        \hline
        \textbf{Trimestre} & \textbf{Contrôle 1} & \textbf{Contrôle 2} & \textbf{Contrôle 3} \\
        \hline
        1\up{er} (sem. 5, 10, 14) & Rappels (pourcentages, évolutions, numération) \newline + Python (conversions) & Compléments analyse + stats desc \newline + Python (box-plot, écart-type) & Probas compléments (indépendance, variance, approximations) \newline + Python (simulation) \\
        \hline
        2\up{e} (sem. 18, 22, 26) & Calcul intégral (primitives, propriétés, théorème) \newline + Python (rectangles) & Intégration (aires, changement de variable, applications) \newline + Python (Monte-Carlo) & Équations différentielles (séparables, linéaires, applications) \newline + Python (Euler) \\
        \hline
        3\up{e} (sem. 30, 34, 38) & Probabilités continues (densités, normale) \newline + Python (simulation) & Estimation, intervalles de confiance \newline + Python (calcul d’IC) & Tests d’hypothèses \newline + Python (tests) \\
        \hline
    \end{tabular}
\end{center}

\subsection{Problèmes basés sur l’apprentissage (PBL) – 2 h chacun}
Trois PBL sont répartis sur l’année (travail en groupes, compte rendu écrit et code Python) :

\begin{enumerate}
    \item \textbf{PBL 1 (semaine 20)} : \emph{Optimisation d’un réacteur}. À partir d’une fonction modélisant le volume d’un réacteur en fonction d’un paramètre, déterminer le volume optimal par calcul intégral et vérifier par optimisation (dérivée).
    \item \textbf{PBL 2 (semaine 26)} : \emph{Modélisation cinétique d’une réaction chimique}. Utiliser une équation différentielle pour modéliser la concentration d’un réactif au cours du temps. Résoudre analytiquement et par la méthode d’Euler. Déterminer le temps de demi-réaction.
    \item \textbf{PBL 3 (semaine 36)} : \emph{Contrôle qualité}. À partir d’un échantillon de mesures, réaliser une estimation ponctuelle et par intervalle de confiance, puis mener un test d’hypothèse pour valider ou non la conformité d’un lot.
\end{enumerate}

\section{Devoirs maison (DM) -- progression, Python et IA}

Un devoir maison est distribué \textbf{toutes les deux semaines}. Chaque DM porte sur l’ensemble du programme traité depuis le début de l’année et comporte un exercice Python et une mission IA.

\begin{center}
    \small
    \rowcolors{2}{gray!8}{white}
    \begin{longtable}{|C{1cm}|C{1cm}|P{6.5cm}|P{6.5cm}|}
        \hline
        \textbf{DM} & \textbf{Sem.} & \textbf{Thèmes évalués (mathématiques + Python)} & \textbf{Mission IA spécifique} \\
        \hline
        \endhead
        \hline
        \multicolumn{4}{r}{Suite à la page suivante} \\
        \endfoot
        \hline
        \endlastfoot

        1 & 3 & Pourcentages et évolutions. Python : calcul de taux moyens. & Demander à l’IA de générer un problème de taux d’évolution sur 3 ans et de le résoudre. \\
        \hline
        2 & 5 & Numération et simplifications. Python : conversions binaire/hexa. & Utiliser l’IA pour vérifier les conversions et expliquer les erreurs. \\
        \hline
        3 & 7 & Fonctions usuelles (log, trigo). Python : tracé de courbes. & Demander à l’IA d’expliquer le domaine de définition de $\ln(x)$ et $\sqrt{x}$. \\
        \hline
        4 & 9 & Optimisation (dérivée quotient). Python : recherche d’extremum. & Utiliser l’IA pour critiquer le tableau de variations obtenu. \\
        \hline
        5 & 11 & Statistique desc (histogrammes). Python : tracé d’histogrammes. & Demander à l’IA d’interpréter un histogramme donné. \\
        \hline
        6 & 13 & Statistique desc (boxplot, écart-type). Python : calcul de paramètres. & Utiliser l’IA pour comparer deux séries à partir de leurs boxplots. \\
        \hline
        7 & 15 & Probabilités (indépendance, variance). Python : simulation. & Demander à l’IA de démontrer la formule de la variance d’une somme de variables indépendantes. \\
        \hline
        8 & 17 & Primitives et intégrales. Python : méthode des rectangles. & Utiliser l’IA pour expliquer la convergence de la méthode des rectangles. \\
        \hline
        9 & 19 & Théorème fondamental. Python : calcul d’intégrales. & Demander à l’IA de démontrer le théorème fondamental sur un exemple. \\
        \hline
        10 & 21 & Aires et volumes (PBL 1). Python : calcul d’aire. & Utiliser l’IA pour proposer une méthode alternative de calcul d’aire. \\
        \hline
        11 & 23 & Applications intégrales (valeur moyenne). Python : calcul de moyenne. & Demander à l’IA de relier la valeur moyenne à une interprétation physique. \\
        \hline
        12 & 25 & Équations différentielles séparables. Python : Euler. & Utiliser l’IA pour comparer la solution exacte et l’approximation d’Euler. \\
        \hline
        13 & 27 & Équations différentielles linéaires. Python : résolution. & Demander à l’IA de trouver une solution particulière pour $y'+2y=e^x$. \\
        \hline
        14 & 29 & Variables aléatoires continues. Python : simulation. & Utiliser l’IA pour expliquer la notion de densité avec un exemple. \\
        \hline
        15 & 31 & Loi normale. Python : calcul de probabilités. & Demander à l’IA de calculer $P(-1<X<1)$ pour $X\sim \mathcal{N}(0,1)$ et d’interpréter. \\
        \hline
        16 & 33 & Échantillonnage et IC. Python : calcul d’intervalles. & Utiliser l’IA pour discuter de l’influence de la taille de l’échantillon sur la largeur de l’IC. \\
        \hline
        17 & 35 & Tests d’hypothèses. Python : tests. & Demander à l’IA d’expliquer la notion de p-valeur. \\
        \hline
        18 & 37 & Synthèse inférence. Python : projet complet. & Utiliser l’IA pour un retour sur le code et proposer des améliorations. \\
        \hline
    \end{longtable}
\end{center}

\section{Conseils pour la mise en œuvre}
\begin{itemize}
    \item \textbf{Rituels} : 5 à 10 minutes de calcul mental et d’extraits de code Python (lecture, débogage) par séance.
    \item \textbf{Travail en îlots} : favoriser les échanges lors des recherches et des PBL.
    \item \textbf{Numérique} : Python régulier (une séance sur deux) avec des modules simples (\texttt{math}, \texttt{matplotlib}, \texttt{scipy.stats}).
    \item \textbf{IA éthique} : apprendre aux élèves à formuler des requêtes précises et à évaluer les réponses de l’IA de manière critique.
    \item \textbf{Différenciation} : proposer des exercices progressifs, des approfondissements pour les plus à l’aise.
\end{itemize}

\section*{Conclusion}
Cette programmation couvre l’intégralité du programme de 2\up{me} année de STS BIO AC, tout en consolidant les lacunes identifiées en 1\up{re} année (notions marquées en rouge). La progressivité spiralaire, l’intégration systématique de Python et de l’IA, ainsi que les évaluations variées (formatives, sommatives, PBL, DM) permettent de développer les compétences de modélisation, de raisonnement et de codage indispensables à la poursuite d’études et à la pratique professionnelle en laboratoire.

\end{document}